\small\noindent\textbf{Abstract}\quad
This simulation-only computational study compares PID, discrete linear-quadratic regulation (LQR), and linear-quadratic integral (LQI) control for an idealized, well-mixed aqueous thermal process. Exact zero-order-hold propagation, independently checked Riccati solutions, equal-budget tuning, a $\pm30\%$ uncertainty domain, constant thermal loads, sensor noise, actuator saturation, and a delayed-sensor plant are evaluated in Rust. Equal-budget tuning changes nominal integral absolute error (IAE) ranking from LQI--PID--LQR to LQR--LQI--PID, but the tuned LQR becomes strongly saturation-dominated under a \SI{20}{s} sensor lag and \SI{10}{s} delay. Integral action removes feasible constant-load offset; steady-noise differences are statistically detectable but practically small. Validated interval bounds certify unsaturated LQR stability, and subdivided interval Jury inequalities certify LQI stability throughout the declared continuous parameter box. PID remains supported by a dense grid rather than a formal interval proof. The contribution is a reproducible characterization of conditional performance trade-offs, not a novel controller, physical experiment, or universal robustness claim.

\par\vspace{0.7em}\noindent\textbf{Keywords}\quad PID control; linear-quadratic control; integral action; parametric uncertainty; thermal process; computational reproducibility.

